\documentclass[12pt]{article}

\usepackage[margin=3cm]{geometry}
\usepackage{amsmath, amssymb, amsthm}
\usepackage{tikz}
\usetikzlibrary{positioning}
\usepackage{xcolor}
\usepackage[inline]{enumitem}
\usepackage{tcolorbox}
\tcbuselibrary{breakable}
\usepackage{fancyhdr}
\usepackage{url}

% Initialize line spacing to 1.5
\linespread{1.5}

% Convenience commands for colored boxes
\newcommand{\rbox}[1]{\fcolorbox{red}{white}{#1}}
\newcommand{\bbox}[1]{\fcolorbox{blue}{white}{#1}}

% Convenience command for hidden boxed content
\newcommand{\hide}[1]{\boxed{\phantom{#1}}}

% Simple exercises environment using basic enumerate
\newcounter{exercisenum}
\newenvironment{exercises}{%
  \begin{tcolorbox}[colback=white,colframe=black,arc=0pt,breakable]
%   \noindent\textbf{Exercises}\par
  \begin{enumerate}[label=\arabic*),leftmargin=1cm,resume=exercises]
  \setcounter{enumi}{\value{exercisenum}}
}{%
  \setcounter{exercisenum}{\value{enumi}}
  \end{enumerate}
  \end{tcolorbox}
}

% Setup footer for first page only
\fancypagestyle{firstpage}{
  \fancyhf{}
  \fancyfoot[L]{\small LYRICS}
  \fancyfoot[R]{\small \quad Ewan Short \quad \today}
  \fancyfoot[C]{\small\thepage}
  \renewcommand{\headrulewidth}{0pt}
  \renewcommand{\footrulewidth}{0pt}
}

\begin{document}

\thispagestyle{firstpage}

% \section{Fractions}
\begin{itemize}
\item
Song lyrics are usually written how they are sung, without worrying about spelling or punctuation. For example, here is part of Olivia Foai'i's song \textit{Party for 1}.
\begin{quotation}
\noindent I'll be messy bunnin' \\
While your party's pumpin' \\
I'm in a t-shirt dancin' on my own \\
Lookin' like risky business \\
Cookin' in the kitchen \\
Don't try to reach me \\
Can't come to the phone
\end{quotation}
\end{itemize}

\begin{exercises}
\item    
Write out some of the lyrics to a song you like. You can find most lyrics on \url{genius.com}. Please choose lyrics that are G or PG rated! \label{Ex:lyrics} If you don't like music, do a TV theme-song, hymn, nursery rhyme etc.~instead.
\item
Why did you choose these lyrics? Are they important to you or relatable in some way? Do they go well with the music?
\end{exercises}

\begin{itemize}
\item
Lyrics are written in lines to emphasize the flow of the music. To see the difference, let's rewrite \textit{Party for 1} as a paragraph. We'll also put the ``proper'' spelling and punctuation back in. 
\begin{quotation}
\noindent I'll be messy bunning while your party's pumping. I'm in a t-shirt, dancing on my own. Looking like risky business, cooking in the kitchen. Don't try to reach me, I can't come to the phone.
\end{quotation}
\item
Olivia Foa'i invented a new word ``bunning" in \textit{Party for 1}. Notice messy is an \underbar{adjective} and bun is a \underbar{noun}, but adding the ``ing'' turns ``bunning" into a \underbar{verb}.
\end{itemize}

\begin{exercises}
\item
Rewrite your lyrics from \ref{Ex:lyrics} as a paragraph using ``proper'' spelling and punctuation.
\item
Why do you think singers invent new words in their music? What mental image does ``messy bunnin" evoke? Did the artist invent any new words in your lyrics from \ref{Ex:lyrics}?
\end{exercises}

\begin{itemize}
\item
A good way to learn creative writing is to imitate the style of people you respect. Here are some lyrics I've written in the style of Olivia Foai'i's \textit{Party for 1}.
\begin{quotation}
\noindent I'll be gluin' Warhammer \\
While you're rockin' bangers \\
I'm in my PJs paintin' river trolls \\
Lookin' like agin' Chewbacca \\
I'll watch A New Hope after \\
Don't try to text me \\
While yellow text scrolls
\end{quotation}
\end{itemize}

\begin{exercises}
\item
Can you write your own lyrics in the style of the song you chose in \ref{Ex:lyrics}?
\end{exercises}

\end{document}